\section{Gemini: ML coding y long horizon}

Después de incorporarse a Google DeepMind, el trabajo de Yao Shunyu (姚顺宇) se centró en ML coding y long horizon. ML coding apunta a un objetivo más radical: que la IA participe, o incluso complete parcialmente, el proceso de «entrenar IA», lo cual incluye elegir los datos, diseñar las señales de retroalimentación y construir el pipeline de entrenamiento y evaluación. Long horizon, por su parte, se refiere a que el modelo mantenga el objetivo, el estado y el contexto de forma consistente en tareas muy largas.

\begin{figure}[H]
\centering
\includegraphics[width=0.88\textwidth]{frame-gemini-training.jpg}
\caption{Capítulo de Gemini: se discuten ML coding, long horizon y el entrenamiento finito con un uso casi infinito.\protect\footnotemark}
\end{figure}
\footnotetext{Intervalo de tiempo en el video: 02:51:43--02:52:15. Este fotograma proviene de la parte intermedia de la discusión sobre el enfoque de trabajo en Gemini.}

\begin{knowledgebox}{Lectura de la figura: lo importante del capítulo sobre long horizon no es la imagen, sino la estructura de la tarea}
Este fotograma corresponde a la discusión sobre ML coding y long horizon. El lector debe poner atención en la tabla que aparece más adelante: compresión de estado, descomposición de tareas, retroalimentación del entorno y memoria/recuperación. La imagen ofrece un punto de tiempo rastreable; la tabla es la que ofrece el desglose conceptual.
\end{knowledgebox}


\begin{importantbox}{train with finite, use as infinite}
Un lema clave de la entrevista es: entrenar con un context finito, pero usarlo de forma casi infinita. Esto no necesariamente significa que una sola muestra de entrenamiento se vuelva infinitamente larga, sino que el modelo aprenda, dentro de un contexto finito, a comprimir, seleccionar, olvidar, recuperar y continuar, de modo que pueda completar tareas que superan por mucho la longitud de la ventana de entrenamiento.
\end{importantbox}

\subsection{Los problemas técnicos de long horizon}

La dificultad de las tareas de largo plazo no es solo la longitud del contexto. Los problemas reales incluyen: qué información se debe conservar, qué información se debe descartar, cómo recuperar el objetivo a partir de estados pasados, cómo evitar que los errores tempranos se acumulen como una bola de nieve, cómo dividir una tarea de largo plazo en subtareas verificables, y cómo evaluar el «mantenimiento del objetivo» en lugar de evaluar solo la respuesta de un único paso.

\noindent\begin{tabularx}{\textwidth}{p{0.22\textwidth}p{0.33\textwidth}X}
\toprule
Problema & Mecanismo & Modo de falla \\
\midrule
Compresión de estado & Resumir las interacciones históricas en una memoria utilizable & Se pierden restricciones clave, lo que provoca una desviación posterior. \\
Descomposición de tareas & Dividir el objetivo de largo plazo en pasos de corto plazo verificables & Las subtareas son correctas, pero el objetivo general es incorrecto. \\
Retroalimentación del entorno & Corregir el comportamiento con resultados de ejecución, pruebas y retroalimentación del usuario & La retroalimentación es escasa o ruidosa, y el modelo aprende algo incorrecto. \\
Memoria/recuperación & Recuperar información relevante desde un almacenamiento externo & Se recupera información similar pero incorrecta. \\
\bottomrule
\end{tabularx}

\subsection{ML coding: que la IA entrene a la IA}

El ML coding al que se refiere Yao Shunyu no consiste solo en que el modelo escriba código de negocio ordinario, sino en que el modelo participe en el sistema de entrenamiento mismo. Esto implica ocuparse de la selección de datos, las señales de retroalimentación, la infraestructura de entrenamiento, el registro de experimentos, la evaluación y el failure analysis. Si esta dirección se logra, cambiará la forma de trabajar de los investigadores: en lugar de escribir a mano toda la lógica de los experimentos, los investigadores pasarían a diseñar objetivos, restricciones, verificaciones y evaluaciones.

\begin{warningbox}{El autoentrenamiento de la IA no es un instituto de conducción autónoma}
Que la IA participe en el entrenamiento de la IA no equivale a soltarle todo el control. La pregunta de entrevista que Yao Shunyu menciona más adelante también lo ilustra: si el candidato le deja todo a la IA sin poder explicar qué hizo la IA, eso se descubre en la sesión de discusión. El valor del investigador del futuro está en colaborar con la IA, entender lo que produce y verificar los resultados, no en confiar a ciegas en la salida.
\end{warningbox}

\subsection{Resumen de la sección}

El núcleo del capítulo de Gemini son dos direcciones: ML coding y long horizon. La primera hace que la IA participe en el proceso de entrenamiento del modelo; la segunda hace que el modelo mantenga el objetivo y el estado en tareas largas. Ambas apuntan en conjunto hacia un flujo de trabajo de IA más autónomo, pero ambas dependen de un diseño cuidadoso de los datos, la retroalimentación y la evaluación.
